\section{Introduction}\label{sec:intro}

A personalised anomaly detector compares each new transaction with a model of one person's own past behaviour. The appeal is plain: a 900-unit purchase is unremarkable for one account holder and extraordinary for another, and a single population-wide rule serves both badly.

The approach carries an equally plain dependency. The personal model has to be estimated from the person's own transactions, and when an account opens there are none. Somewhere between zero transactions and a lifetime of them, the personal baseline becomes worth using. Production systems settle the question with a cutoff; one patent family describes filtering out ``unseasoned users, i.e.\ users with insufficient history'' \cite{faigon2021anomaly}. The cutoff is treated as a known engineering parameter, yet we found no published measurement of where it lies or whether it exists (\cref{sec:related}).

We measure it. For each account we hold a 60-transaction evaluation window fixed and train on the $n$ transactions immediately preceding it, with $n$ from 1 to 250 on real data and to 500 on synthetic data, so that every comparison across history sizes is paired at the account level. Anomalies are defined by provenance: a held-out transaction is replaced by a real transaction from a different account, so no detector's score can define its own labels. We run six unsupervised detectors on 300 accounts of the PKDD'99 Berka bank data, on a controlled simulation of 240 individuals whose variability we set directly, and on 300 card entities of IBM's synthetic TabFormer data, which carries merchant, category and geographic channels that Berka lacks.

The shape of the answer is an early plateau. A shrinkage-regularised Mahalanobis detector reaches 90\% of its best precision-recall AUC (PR-AUC, 0.657 against a 0.10 prevalence baseline) at 25 personal transactions, and the threshold is 25 at each of five independent cohort draws. The single number conceals three dependencies that the rest of the paper quantifies: the threshold moves with the detector and the dataset, the usable alert threshold moves with the calibration rule, and below the sufficiency point a population model beats the person's own history.

\paragraph*{Contributions}
\begin{enumerate}
\item \textbf{A measured sufficiency point.} Data-efficiency curves for per-user anomaly detection under a provenance-labelled, paired protocol. On Berka, Mahalanobis distance reaches 90\% of its best PR-AUC at 25 personal transactions and gains nothing measurable beyond 100; the same early-plateau shape appears on TabFormer, with a threshold of 250 (\cref{sec:curves,sec:tabformer}).
\item \textbf{A crossover and its remedy.} A population model outperforms the personal model below about 10 transactions for Mahalanobis and 50 for $k$-NN on Berka. A history-weighted shrinkage hybrid is never meaningfully worse than either parent at any size on Berka and removes the cold-start penalty without a measurable cost at large $n$ (\cref{sec:arms}).
\item \textbf{A calibration effect that changes the operating point.} Calibrating the alert threshold on the model's own training scores makes Mahalanobis flag 96\% of transactions at $n=10$; leave-one-out (LOO) calibration makes it flag 1.2\% there and brings the false-positive rate within 1.6 to 3 points of the nominal 5\% from $n=25$. The effect holds across six detectors and five seeds, and it does not always push in the same direction (\cref{sec:calibration}).
\item \textbf{The limits of transfer.} Detector rankings under scarcity reverse between datasets (Isolation Forest is weakest on Berka and strongest at small $n$ on TabFormer); behavioural variability predicts the requirement under controlled manipulation (Spearman $\rho=0.48$ to $0.71$) but not across real accounts; and a direct drift measurement rules out drift as the explanation for the negative marginal gain from 100 to 250 transactions (\cref{sec:variability,sec:recency,sec:tabformer}).
\end{enumerate}

\Cref{sec:related} positions the work, \cref{sec:setting} defines the setting and data, \cref{sec:protocol} the protocol, and \cref{sec:results} reports results. \Cref{sec:discussion} draws out what the evidence establishes and \cref{sec:limits,sec:ethics} give limitations and ethical considerations.
